\documentclass[11pt]{article}
\usepackage[a4paper,margin=25mm]{geometry}
\usepackage{amsmath,amssymb,amsthm,booktabs}
\setlength{\emergencystretch}{2em}
\newcommand{\gen}[1]{\langle#1\rangle}
\newcommand{\Sym}{\operatorname{Sym}}
\newcommand{\Alt}{\operatorname{Alt}}
\newcommand{\eps}{\varepsilon}
\newtheorem{theorem}{Theorem}
\newtheorem{lemma}[theorem]{Lemma}
\title{Signed permutation groups at the ends of alternating group constructions}
\author{Human Compression project}
\date{Author edition; distinct review pending}
\begin{document}\maketitle

This note supplies uniform proofs for the higher odd-degree families in
Fernandes--Leemans, arXiv:1810.12450v2, Theorems~4.8--4.10. The literal
permutation families belong to that source. We describe their actual end
subgroups, including the relations between their separate orbits. The
argument uses the standard end criterion for string C-groups, the acquired
restriction-kernel theorem, and the supported alternating-group lemma from
our reviewed companion note \emph{Finite end groups for alternating group
constructions}. The earlier maximal-rank family, explicitly attributed
to [14] in source Theorem~4.7, is retained as an input, not counted as a
new result of this note.

All signs below take values in $\mathbb F_2$. We regard symmetric and
alternating groups on named sets as supported permutation groups, fixing
every other point. A C-marking has the string commutations and the
intersection property for every pair of subsets of its marks.

\section{Two elementary transfer principles}
The acquired restriction-kernel theorem says the following. For a C-marking
$(h_1,\ldots,h_d)$, suppose only $h_1$ moves a point $a$. Put
$P_j=\gen{h_1,\ldots,h_j}$ and $X_j=a^{P_j}$, with $P_0=1$.
Adjoining the pure
transposition $(c\ a)$ at a new point preserves C precisely when
\begin{equation}\label{leaf}
 P_j\cap\Sym(X_{j-1})\le P_{j-1}\quad(1\le j\le d).
\end{equation}
The first test is trivial. The group after adjoining the transposition is
\begin{equation}\label{factor}
 \Sym(X_d\cup\{c\})\times P_d^{\Omega\setminus X_d}.
\end{equation}
The second factor is the actual restriction image. It follows by
conjugating $(c\ a)$ to all star transpositions, then cancelling the
action on $X_d$. Once the first mark is pure, another pure leaf may
always be attached. We also use the proved forward endpoint twist:
in a direct product with a central involution $z$, replacing an endpoint
mark $h$ by $hz$ preserves C. No converse to this assertion is used.

\begin{lemma}[The quotient criterion]\label{quotient}
Suppose a homomorphism maps a string marking onto a C-marking, and is
injective on one endpoint deletion of the original marking. Then the
original marking is C.
\end{lemma}
\begin{proof}
Reverse the ordering so that injection holds on deletion of the first
mark. Induct on the length. That endpoint deletion is isomorphic, as a
marked group, to its image and hence is C. The other endpoint deletion
is C by induction, since its own deletion of the first mark is inside
the subgroup on which the homomorphism is injective. An element in the
intersection of the two endpoint groups maps into the middle group
by the quotient's intersection property. Choose a word in the middle
marks with that image. Both this word and the original element lie in
the injective endpoint group, so they are equal. The end criterion
proves the assertion. The length-one start is immediate. This is the
usual McMullen--Schulte quotient criterion, Proposition~2E17.
\end{proof}

\section{Two markings of the even signed permutation group}
On $\{1^+,1^-,\ldots,t^+,t^-\}$ let $s_i$ interchange the two
coordinate pairs $i,i+1$, preserving their signs. Let $v_i$ interchange
$i^+,i^-$. Thus
\[
 D_t=\{v_1^{e_1}\cdots v_t^{e_t}:\textstyle\sum e_i=0\}
       \rtimes S_t.
\]
The symbol $D_t$ in this note denotes this concrete permutation group,
of order $2^{t-1}t!$, rather than a dihedral group.

There are two strings. The first, for $t\ge5$, is
\begin{equation}\label{two}
 s_1,\ldots,s_{t-4},\quad T=s_{t-3}v_{t-1}v_t,
       \quad s_{t-2},\quad s_{t-1}.
\end{equation}
The second, for $t\ge7$, is
\begin{equation}\label{four}
 s_1,\ldots,s_{t-6},\quad
 T=s_{t-5}v_{t-3}v_{t-2}v_{t-1}v_t,
 \quad s_{t-4},\quad s_{t-3}s_{t-1},\quad s_{t-2}.
\end{equation}
Empty initial lists have their usual meaning. In either case write
$f$ for the displayed product of $v_i$ in $T$, and $p$ for its
unsigned index: $p=t-3$ or $p=t-5$. Directly,
\begin{equation}\label{cube}
 (s_{p-1}T)^3=f.
\end{equation}
Indeed the unsigned product is a three-cycle, and $f$ is disjoint
from its three coordinates. Consequently the prefix ending in $T$
contains $f$ and every unsigned $s_i$ for $i\le p$.

\begin{lemma}\label{caps}
Both strings are C-markings of $D_t$. For fresh points $c,d$, one
may prepend $(c\ 1^+)$, then prepend $(d\ c)$.
After the two attachments the group is
$S_{2t+2}$. In \eqref{two}, deletion of the last mark after these
attachments gives $S_{2t}\times C_2$ on a set of size $2t$ and the
last signed pair. In \eqref{four}, it gives
$S_{2t-2}\times V_4$, where $V_4$ is the regular four-point group
on the last two signed pairs.
\end{lemma}
\begin{proof}
We first record the prefix groups before proving C. The prefix ending
in $T$ is
\begin{equation}\label{N}
 N_t=\operatorname{diag}S_{p+1}\times\gen f.
\end{equation}
Here the symmetric group acts on the first $p+1$ signed pairs without
sign changes; $f$ has disjoint support. On adjoining $s_{p+1}$, the
prefix becomes
\begin{equation}\label{M}
 M_t=\{v_1^{e_1}\cdots v_{p+2}^{e_{p+2}}
       (v_{p+3}\cdots v_t)^{\sum e_i}:e_i\in\mathbb F_2\}
       \rtimes S_{p+2}.
\end{equation}
Conjugates of $f$ under the unsigned $S_{p+2}$ give each vector
$v_i v_{p+3}\cdots v_t$, proving the equality. This is the usual
signed symmetric group on the first $p+2$ coordinates, with every
remaining sign recording the total number of their flips.

For \eqref{two}, $M_t$ is the penultimate prefix. The final unsigned
bridge gives $S_t$. The conjugates of $f$ give all even sign changes:
subtracting two conjugates which differ at just two coordinates gives
$v_i v_j$, and these vectors span the even subspace. Thus the full
group is $D_t$.

For \eqref{four} there is one more intermediate prefix. Adjoin
$b=s_{t-3}s_{t-1}$ to \eqref{M}. Its group is
\begin{equation}\label{L}
 L_t=D_{t-2}\times
       \gen{v_{t-1}v_t,s_{t-1}}=D_{t-2}\times V_4.
\end{equation}
To see that the unsigned factors separate, take the preceding unsigned
transposition $a=s_{t-4}$. The cube $(ab)^3$ is $s_{t-1}$;
hence $s_{t-3}$ is also available. We have unsigned
$S_{t-2}\times S_{\{t-1,t\}}$. Its conjugates of $f$ are
$v_i v_j v_{t-1}v_t$ with $i,j\le t-2$. The product for the
three pairs of a triangle gives $v_{t-1}v_t$ alone. Removing this
factor gives all even flips on the first $t-2$ coordinates. These
are exactly the independent factors in \eqref{L}. The final bridge
$s_{t-2}$ gives unsigned $S_t$, and the same even-vector argument
gives the full $D_t$.

All distant marks commute by inspection. For C, the two finite starts
are \eqref{two} at $t=5$ and \eqref{four} at $t=7$. The certificate
checks all $16^2=256$ and $32^2=1024$ ordered pairs of subset
subgroups, respectively, as actual permutation groups. These are
finite premises; their extension is the following induction.

The prefix $N_t$ is an endpoint central twist of the ordinary
adjacent-transposition simplex and is C. At every subsequent prefix,
deleting its first mark gives the corresponding prefix with $t-1$
coordinates, fixing the first signed pair pointwise. Deleting its
last mark gives the preceding prefix in \eqref{N}--\eqref{L}.
The intersection is exactly the latter prefix's pointwise stabilizer
of the first signed pair. Equations \eqref{N}--\eqref{L} show that
this stabilizer is the preceding prefix with $t-1$ coordinates:
the unsigned permutation fixes coordinate one, and its sign there
is zero; the same parity equation remains on the other coordinates.
This applies successively to $M_t$, then $L_t$ when present, then
$D_t$. The finite starts contain all their own prefix submarkings,
so induction beginning at $t=6$ or $t=8$ proves every case.

For the first pure leaf use \eqref{leaf}. Until $M_t$ is reached,
the preceding root orbit consists only of positive signed points.
A signed permutation supported on such a set is the identity, since
it commutes with the involution interchanging every positive and
negative point. At the subsequent stages the preceding root orbit
contains both signs of an initial coordinate segment. In
\eqref{two}, the only remaining kernel is $D_{t-1}$, supported
on the first $t-1$ coordinates, and it is contained in $M_t$.
In \eqref{four}, the two remaining kernels are $D_{t-3}$ inside
$M_t$ and $D_{t-2}$ inside $L_t$. These are exactly the supported
subgroups obtained from the formulas. Thus all kernel inclusions
hold. A second pure leaf is permitted by the pure-leaf corollary.
The full $D_t$ is transitive on its $2t$ signed points, so
\eqref{factor} twice gives $S_{2t+2}$. For the last deletions,
$M_t$ in \eqref{two} has root orbit of size $2t-2$ and complement
image $C_2$; $L_t$ in \eqref{four} has root orbit of size $2t-4$
and complement image $V_4$. Applying \eqref{factor} twice gives
the asserted actual direct products.
\end{proof}

\section{The left path and its two parity relations}
Let $k\ge1$, $m=4k+6$, and $L=m-1$. On $x_1,\ldots,x_m$ put
\begin{align*}
 a_0&=(x_1\ x_2)(x_3\ x_4)\cdots(x_{L-2}\ x_{L-1}),\\
 a_1&=(x_2\ x_3)(x_4\ x_5)\cdots(x_{L-1}\ x_L),\\
 a_2&=(x_3\ x_4)(x_L\ x_m).
\end{align*}
Add two rows by defining, for $j\ge3$,
\begin{equation}\label{bridges}
 a_j=(x_{m+j-3}\ x_{m+j-2})(y_{j-2}\ y_{j-1}).
\end{equation}
These are ordinary bridges, without the final vertical modifications
which will be introduced below. Write $H_j=\gen{a_0,\ldots,a_j}$
and $K_j=\gen{a_1,\ldots,a_j}$.

\begin{lemma}\label{paths}
The marking of $H_2$ is C and $H_2=\Alt(\{x_1,\ldots,x_m\})$.
For $j\ge3$, its extension by \eqref{bridges} is C, and
\begin{equation}\label{P}
 H_j=\{(\sigma,\pi)\in S_{m+j-2}\times S_{j-1}:
                         \eps(\sigma)=\eps(\pi)\},
\end{equation}
on the two indicated rows. The groups $K_j$, for $j\ge2$, are
C and have the actual form
\begin{equation}\label{K}
 \{(\sigma,\pi,d)\in S_{j+1}\times S_{j-1}\times D:
              \eps(\sigma)+\eps(\pi)=\chi(d)\}.
\end{equation}
Here $D=\gen{a,b}\cong\mathrm{Dih}(8)$ acts on
$\{x_2,x_3,x_4,x_5\}$, with
$a=(x_2\ x_3)(x_4\ x_5)$, $b=(x_3\ x_4)$, and
$\chi(a)=\chi(b)=1$. The first symmetric factor acts on
$\{x_{m-2},\ldots,x_{m+j-2}\}$, and the second on the bottom
row. There are $2k-1$ further two-point orbits, all interchanged
by $a_1$ and fixed by the other marks. Their common action is
determined by the character of $D$ taking $a$ to one and $b$ to zero.
\end{lemma}
\begin{proof}
The element $(a_1a_2)^4$ is a pure three-cycle on
$\{x_{L-1},x_L,x_m\}$. The other nontrivial components of
$a_1a_2$ have lengths four or two. Each of $a_0,a_1,a_2$ fixes
a point of this three-element support, and their graph on the top
row is connected. The acquired supported-alternating-group lemma
therefore gives $H_2=A_m$. All three marks are even. For its
intersection property, $\gen{a_0,a_1}$ is dihedral of order $2L$
and fixes $x_m$. The group $\gen{a_1,a_2}$ is dihedral of order
24; its orbit at $x_m$ has three points, so that point stabilizer
has order eight. Since $L$ is odd, the intersection has order
dividing $\gcd(2L,8)=2$. It contains $a_1$, and hence is exactly
$\gen{a_1}$. The rank-three end criterion applies.

Consider next the two right paths in $K_j$, temporarily omitting
$D$ and the repeated pairs. Their marking consists of two successive
pure leaves at the beginning of a diagonal adjacent-transposition
path on $j-1$ paired coordinates. That diagonal path is C. Its
restriction kernels on an all-positive root orbit are trivial, so
the first leaf satisfies \eqref{leaf}; the second is pure.
Formula \eqref{factor} shows that the resulting quotient is
$S_{j+1}\times S_{j-1}$, with its stated C-marking. This includes
$j=2$, when it is just the two adjacent transpositions of $S_3$.

The two initial quotient marks are odd in total, and every subsequent
paired bridge is even. This gives the relation in \eqref{K}.
The cube $(a_1a_2)^3$ has trivial quotient action and induces
$(ab)^3$ on $D$, generating the cyclic subgroup $\ker\chi$ of
order four. The repeated-pair actions are already determined by
$D$. Thus the kernel is all of $\ker\chi$, and \eqref{K} is
an equality. On omission of $a_1$, the $D$-action reduces to
$\gen b$ and is determined by the quotient's total sign.
Projection is injective there. Lemma~\ref{quotient} proves C for
$K_j$.

To prove \eqref{P}, the normal closure of the supported $A_m$
under the bridges contains the alternating group on the entire
upper row. This follows from the supported-alternating-group lemma
on that row: all conjugates remain supported there, the upper
graph is connected, and each mark fixes a point of the initial
set of $m$ points. The lower projection is $S_{j-1}$. Every even
lower permutation is represented by an even-length word in its
adjacent transpositions; the corresponding bridge word also acts
evenly on the upper row, and that action can be cancelled using
its supported alternating group. We therefore have both supported
alternating factors, together with a bridge odd in both factors.
This proves precisely \eqref{P}.

Finally $H_{j-1}$ fixes the newly added last point on each row.
Inside \eqref{K}, imposing these two fixed points replaces
$S_{j+1}\times S_{j-1}$ by $S_j\times S_{j-2}$, with the same
$D$ relation; the result is exactly $K_{j-1}$. Consequently
$H_{j-1}\cap K_j=K_{j-1}$. All $K_j$ are C and $H_2$ is C,
so induction with the end criterion proves C for every $H_j$.
\end{proof}

\section{Higher odd ranks}
Let $r\ge7$ be odd, $k\ge1$, and $n=2r+1+4k$. Use the preceding
marks through $a_{r-1}$ on top row
$x_1,\ldots,x_{m+r-3}$ and bottom row $y_1,\ldots,y_{r-2}$,
except that the mark $a_{r-3}$ is multiplied by
\begin{equation}\label{vertical2}
 f=(x_{m+r-4}\ y_{r-3})(x_{m+r-3}\ y_{r-2}).
\end{equation}
Call the resulting marks $\rho_0,\ldots,\rho_{r-1}$ and their
group $G$. These are the literal source Theorem~4.8 diagrams.
The two last vertical pairs will be denoted $Z'$ and $Z$, in
that order. All marks are even involutions and distant marks
commute. Write $G_I=\gen{\rho_j:j\notin I}$.

The right end $G_0$ has its four-point $D$-orbit, repeated pairs,
and a large orbit $R$ of size $N=2r-2$. On $R$, its marking is
\eqref{two} with $t=r-2$, preceded by two pure leaves. Thus
Lemma~\ref{caps} gives the quotient $S_R$ with C. The same
cube and injectivity argument as in \eqref{K} gives
\begin{equation}\label{G0odd}
 G_0=\{(g,d)\in S_R\times D:\eps(g)=\chi(d)\},
\end{equation}
including the determined repeated-pair actions, and proves C.
Deleting its last mark gives, from the last-deletion part of
Lemma~\ref{caps},
\begin{equation}\label{middleodd}
 G_{0,r-1}=
 \{(g,z,d)\in S_{R\setminus Z}\times S_Z\times D:
                  \eps(g)+\eps(z)=\chi(d)\}.
\end{equation}
This is also C as a submarking of $G_0$.

We prove C for the other end $G_{r-1}$. Its deletion of the last
remaining mark is $G_{r-2,r-1}$. By Lemma~\ref{paths}, the
ordinary prefix through $a_{r-3}$ is a C-marking of the parity
group on
\[
 U=\{x_1,\ldots,x_{m+r-5}\},\qquad
 V=\{y_1,\ldots,y_{r-4}\}.
\]
Multiplying its endpoint by the disjoint central $f$ preserves C.
The adjacent preceding bond is three, so the cube separates $f$.
Consequently this end is exactly
\begin{equation}\label{Htwist}
 G_{r-2,r-1}=\{(u,v)\in S_U\times S_V:
                             \eps(u)=\eps(v)\}\times\gen f.
\end{equation}
Its intersection with \eqref{middleodd} can now be calculated
without a maximal-subgroup assertion. Put
$U'=U\cap R$, of size $r-2$. An element of the intersection
preserves $U',V,Z',Z$, and the actions on $Z',Z$ are the same
interchange, say with exponent $e$. On the rest of $U$, the
four-point $D$-action has sign equal to the $b$-character, while
the odd number $2k-1$ of repeated pairs contributes the
$a$-character. Their sum is $\chi(d)$. Thus \eqref{Htwist}
imposes exactly
\[
 \eps(u|_{U'})+\eps(v)=\chi(d).
\]
The relation in \eqref{middleodd} is identical, since the two
pair signs $e+e$ cancel. Therefore the intersection is
\begin{equation}\label{oddintersection}
 \{(u',v,d)\in S_{r-2}\times S_{r-4}\times D:
              \eps(u')+\eps(v)=\chi(d)\}\times\gen f.
\end{equation}
By \eqref{K} and the same separating cube, this is exactly
$G_{0,r-2,r-1}$. Both ends of $G_{r-1}$ are C, so its
intersection property follows.

Finally $G_{r-1}$ preserves the last pair $Z$. The full setwise
stabilizer of $Z$ in \eqref{G0odd} is precisely
\eqref{middleodd}. Hence
$G_0\cap G_{r-1}=G_{0,r-1}$, proving C for $G$.
The left three marks contain the supported $A_m$ proved in
Lemma~\ref{paths}. The full graph is connected, and every mark
fixes a point of that $m$-element set. The acquired alternating
support lemma gives $G=A_n$.

The adjacent orders are
\begin{equation}\label{oddtype}
 [4k+5,12,6,3^{r-7},6,6,3].
\end{equation}
The first is the length of the initial alternating path; the next
is the least common multiple of four and three on the two
nontrivial components of $a_1a_2$. The third is a three-cycle
with disjoint transpositions. The unchanged paired bridges give
three, and the three final products involve only the four last
paired coordinates. Their orders $6,6,3$ follow either directly
from \eqref{two} or from its finite start. Thus these are the
source's exact types as well as its rank-existence assertions.

\section{Higher even ranks}
Let now $r\ge8$ be even, $k\ge0$, $m=4k+6$, and
$n=2r+3+4k$. Keep the initial marks and bridges, extend the top
row to $x_{m+r-2}$ and the bottom to $y_{r-1}$, and make these
two modifications:
\begin{align}
 \rho_{r-2}&=a_{r-2}
       (x_{m+r-3}\ x_{m+r-2})(y_{r-2}\ y_{r-1}),\label{extramatch}\\
 \rho_{r-4}&=a_{r-4}
       \prod_{i=0}^{3}(x_{m+r-2-i}\ y_{r-1-i}).\label{vertical4}
\end{align}
All other marks are the corresponding $a_j$. These are the
source Theorems~4.9 and 4.10 diagrams. They again give even
involutions with the string commutations.

First omit $\rho_{r-1}$. There is a main component $M$ of size
$n-4$ on which the marking is exactly the odd-rank construction
of rank $r-1$, and a four-point component $Q$. On $Q$, only
$\rho_{r-4}$ and $\rho_{r-2}$ act, as the two independent
double transpositions $v,h$ of a square. The main marking is C:
for $k\ge1$ this was just proved, and for $k=0$ it is the
retained earlier maximal-rank seed of source Theorem~4.7.

Projection of $G_{r-1}$ onto $M$ is injective on the endpoint
deletion omitting $\rho_{r-2}$. Indeed $h$ is then absent,
and its action on $Q$ is $v^e$, where $e$ is the parity of the
number of occurrences of $\rho_{r-4}$. On the main component,
the last vertical pair is preserved by that endpoint deletion;
exactly $\rho_{r-4}$ interchanges it. Hence $e$ is determined
by the main permutation. A word acting trivially there has
$e=0$ and acts trivially on $Q$ as well. Lemma~\ref{quotient}
proves C for $G_{r-1}$.

Its actual group is
\begin{equation}\label{evenlast}
 G_{r-1}=\Alt(M)\times V_4.
\end{equation}
For the commutator subgroup has trivial action on $Q$, and its
image on $M$ is $[\Alt(M),\Alt(M)]=\Alt(M)$; here
$|M|\ge15$. Thus the full supported $\Alt(M)$ lies in the
group. Cancelling the main parts of the two indicated marks
then yields $v$ and $h$ separately. This proves \eqref{evenlast}.

For $k\ge1$, the right end has the same $D$-orbit and repeated
pairs as before, and a large orbit $R$ of size $N=2r$. On $R$
the marking is \eqref{four} with $t=r-1$, preceded by two pure
leaves. Exactly as in \eqref{G0odd}, the quotient criterion and
the cyclic kernel of order four give C and
\begin{equation}\label{G0even}
 G_0=\{(g,d)\in S_R\times D:\eps(g)=\chi(d)\}.
\end{equation}
The last deletion, by Lemma~\ref{caps}, is
\begin{equation}\label{middleeven}
 G_{0,r-1}=
 \{(g,d)\in S_{R\setminus Q}\times D:\eps(g)=\chi(d)\}
       \times V_4.
\end{equation}
Every $V_4$-element on $Q$ is even. Equations
\eqref{evenlast} and \eqref{G0even} therefore force their
intersection to be exactly \eqref{middleeven}: preservation of
$M,Q$ leaves the displayed symmetric factor, and the sign on
the rest of $M$ is precisely $\chi(d)$. The end criterion
proves C for the full marking when $k\ge1$.

It remains to justify $G_0$ when $k=0$, where the two left
small orbits join and the formula with $D$ does not apply.
In this case $G_0$ fixes $x_1$ and is transitive on all other
$n-1$ points. The subgroup $G_{0,1}$ has the separate pair
$\{x_3,x_4\}$ and a large orbit of size $n-4$; it fixes
$x_1,x_2$. On that large orbit its marking is one pure leaf
attached to \eqref{four}; the first mark is multiplied by the
interchange of $x_3,x_4$. It is an endpoint central twist, and
its group is
\begin{equation}\label{smallend}
 \{g(x_3\ x_4)^{\eps(g)}:g\in S_{n-4}\}.
\end{equation}
Here only the attached leaf is odd on the large orbit and every
signed-cap mark is even, which proves this precise sign relation.
It contains the supported alternating group on that orbit.
Every mark of $G_0$ fixes a point of the orbit, and its full
graph off $x_1$ is connected. The alternating support lemma
therefore gives
\begin{equation}\label{maxright}
 G_0=\Alt(\Omega\setminus\{x_1\}).
\end{equation}
This generation argument does not presuppose C.

For C, one end of $G_0$ is \eqref{smallend}, already proved C.
Its other end $G_{0,r-1}$ is a submarking of the C-group
$G_{r-1}$. We also need its exact group. On $M$ it is the
right end of the retained odd maximal-rank seed. That right
end is $\Alt(M\setminus\{x_1\})$: the preceding argument
applies verbatim with \eqref{two} in place of \eqref{four}.
The pure-leaf cap gives a large supported alternating group,
and connectivity gives the full even group on the other points.
As in \eqref{evenlast}, taking commutators and then cancelling
the two square actions proves
\begin{equation}\label{k0middle}
 G_{0,r-1}=\Alt(M\setminus\{x_1\})\times V_4.
\end{equation}
Intersect this with \eqref{smallend}. The element fixes $x_2$,
preserves $\{x_3,x_4\}$, and acts on $Q$ in $V_4$. On the
remaining $n-8$ points it is arbitrary symmetric, with its sign
recorded by the interchange of $x_3,x_4$. Thus the intersection
is the sign copy of $S_{n-8}$, times $V_4$. This is exactly
$G_{0,1,r-1}$: use the last-deletion formula of
Lemma~\ref{caps} with one attached leaf and its same endpoint
twist. Both ends are C and their intersection is the middle,
so $G_0$ is C. Finally \eqref{maxright} and \eqref{evenlast}
intersect in \eqref{k0middle}, completing the full C proof.

For every $k\ge0$, the supported $\Alt(M)$ in
\eqref{evenlast}, the connectivity of the full graph, and a
fixed point in $M$ for each mark give $G=A_n$ by the alternating
support lemma. The type is
\[
 [5,5,6,3^{r-8},6,6,6,4]\quad(k=0),\qquad
 [4k+5,12,6,3^{r-8},6,6,6,4]\quad(k\ge1).
\]
The initial products are as above (a five-cycle when $k=0$),
the middle paired bridges give three, and the final four
products are those in \eqref{four}, with orders $6,6,6,4$.

\section{Scope and attribution}
The written proof recovers the exact higher odd-degree families
of source Theorems~4.8--4.10. The parameterization is exhaustive
for their stated ranks: when $n\equiv3\pmod4$, odd $r$ has
$n=2r+1+4k$, and even $r$ has $n=2r+3+4k$. The odd $k=0$
seed remains the attributed prior theorem. This note does not
claim a new construction, a new open-question answer, a Lean
verification, or a numerical compression factor.

The new organization consists of the two concrete signed cap
calculations, the restriction-kernel transfer, and the exact
parity descriptions of the end intersections. In particular,
an abstract group isomorphism or a list of orbit sizes is never
used as a substitute for an equality of permutation subgroups.
The only new finite premises are the two complete base cap
certificates. Sample computations in larger degrees are retained
as diagnostics and are not premises of a uniform assertion.

This component requires distinct review before admission. The
low-rank companion is already reviewed separately. A whole L038 recovery also needs the
existing contraction and degree-twelve interfaces assembled in
one principal scope route.

\begin{thebibliography}{9}
\bibitem{FL} M. E. Fernandes and D. Leemans,
\emph{String C-group representations of alternating groups},
arXiv:1810.12450v2, 4 January 2019; Ars Math. Contemp. 17 (2019),
291--310. Theorems 4.7--4.10 and their literal diagrams.
\bibitem{prior} M. E. Fernandes, D. Leemans and M. Mixer,
\emph{All alternating groups $A_n$ with $n\ge12$ have
polytopes of rank $\lfloor(n-1)/2\rfloor$},
SIAM J. Discrete Math. 26(2) (2012), 482--498.
The maximal-rank construction is retained through its explicit
attribution and statement in source Theorem 4.7.
\bibitem{MS} P. McMullen and E. Schulte,
\emph{Abstract Regular Polytopes}, Cambridge University Press,
2002, Propositions 2E16 and 2E17.
\end{thebibliography}
\end{document}
